\documentclass[11pt]{letter}
\usepackage[T1]{fontenc}
\usepackage[utf8]{inputenc}
\usepackage[margin=1in]{geometry}
\usepackage{lmodern}
\usepackage{hyperref}
\hypersetup{hidelinks}
\setlength{\emergencystretch}{2em}

\signature{Mingyuan Xie\\Zhengxun Wu\\
Qingdao University of Technology\\
Corresponding authors\\
\href{mailto:jhj208436@gmail.com}{jhj208436@gmail.com}\\
\href{mailto:2276815088@qq.com}{2276815088@qq.com}}
\address{Qingdao University of Technology\\Qingdao, Shandong, China}
\date{10 August 2026}

\begin{document}

\begin{letter}{Editor-in-Chief\\\emph{Computing}}

\opening{Dear Editor-in-Chief,}

We submit the manuscript entitled
\emph{``Resource-Aware Permission Governance for Dependable Multi-Tenant LLM
Agent Runtimes''} for consideration as a regular research article in
\emph{Computing}.

The manuscript addresses a systems dependability problem at the intersection
of resource admission and permission governance. Under hostile resource
pressure, a shared policy executor can miss its decision deadline; if the
runtime fails open, this availability failure becomes an unsafe admission. We
present Neuron, a Java runtime architecture that reserves permission
capacity and couples a resource-pressure signal to the escalation policy.
The central contribution is the experimentally separable distinction between
capacity isolation and cross-layer coupling; identity propagation,
privilege non-increase, and correlated auditing are supporting mechanisms.

The evaluation uses production Java~21 token-bucket and escalation-policy
code, four predeclared architectures, independent JVM runs on Windows and an
Ubuntu/KVM VPS, deadline and attacker-pressure sweeps, capacity sensitivity,
process-separated permission-service fault injection, API-derived trace
replay, and a post-selection holdout on six previously unused workload
profiles. The frozen data and one-click replication package contain the raw
JSONL observations, environment records, analysis scripts, generated values,
and source code. The process-separated VPS experiment is explicitly a
same-host loopback study; we do not claim WAN or multi-node behavior.

We believe the manuscript fits \emph{Computing}'s interest in dependable,
adaptive, and trustworthy computing systems because it provides a systematic
runtime foundation and a controlled architectural comparison rather than a
feature-only security prototype. The work is original, is not under
consideration elsewhere, and has been approved by all authors. Both authors
are corresponding authors. We have no competing interests and received no
funding for this study.

Thank you for your consideration.

\closing{Sincerely,}

\end{letter}
\end{document}
